Au repos, le cœur d'une élève bat à \qty{\v{FC}}{battements/min} et chaque contraction du ventricule gauche éjecte \qty{\v{VES}}{mL} de sang dans l'aorte. Pendant une course, sa fréquence cardiaque monte à \qty{\v{FCe}}{battements/min} et chaque contraction éjecte alors \qty{\v{VESe}}{mL} de sang.

\begin{enumerate}[label=\alph*)]
    \item Calculez son débit cardiaque au repos, en litres par minute.
    \item Son corps contient environ \qty{5}{L} de sang. Au repos, combien de temps faut-il pour que l'équivalent de tout ce sang passe par le cœur ?
    \item Calculez son débit cardiaque pendant la course. Par quel facteur a-t-il augmenté ? Expliquez pourquoi les muscles ont besoin de ce débit plus élevé pendant l'effort.
\end{enumerate}
